\section{Operators and pair fibers}\label{sec:setup}
This section fixes the operator conventions and expresses both chains through
the switch Laplacian and the pair-fiber projections. All operators act on
real $L^2(\pi)$. Write $A\succeq B$ when $A-B$ is self-adjoint and has
nonnegative quadratic form. Let $\Pi f=\mathbb E_\pi f$ be projection onto
constants.

Each distinct switch neighbor of a graph has transition probability
$1/(12\binom n4)$ under $P$: the four toggled edges determine both the
four-set and the ordered pair of removed and added matchings, so this count
has no proposal multiplicity. If $L$ is the unnormalized Laplacian of the
graph of distinct switch neighbors, then
\begin{equation}\label{eq:normalization}
 Lf(G)=\sum_{G'\sim G}(f(G)-f(G')),\qquad I-P=\frac{L}{12\binom n4}.
\end{equation}

For an unordered vertex pair $a=\{i,j\}$, its \emph{fiber} fixes all edges
away from $a$, the edge $ij$, and each outside vertex's number of neighbors
in $a$. The graphs in the fiber of $G$ are exactly those obtainable from $G$
by a Curveball trade at $a$. Let $E_a$ be conditional uniform averaging on
this fiber and put
\begin{equation}\label{eq:operators}
 h_a=I-E_a,\qquad H=\sum_{a\in\binom{[n]}2}h_a,\qquad
 K=\frac{1}{\binom n2}\sum_a E_a=I-\frac{H}{\binom n2}.
\end{equation}
This is the partition-based heat-bath construction of Dyer, Greenhill and
Ullrich~\cite[Definition~1.1]{DGU14}, so $E_a$ and $h_a$ are orthogonal
projections. Thus $K$ chooses a uniform unordered pair and performs a uniform
undirected Curveball trade~\cite{CBS18}, including the possibility of
retaining the current assignment. No additional holding probability is
imposed. Both kernels are symmetric and stochastic. Also $P$ has nonnegative spectrum by
laziness, and $K$ does because it is an average of orthogonal projections.
Section~\ref{joh:section} describes each fiber as a uniform subset slice, and
Section~\ref{asm:section} combines the estimates for $H$ with
\eqref{eq:normalization} and \eqref{eq:operators}.
